\documentclass{situationist-preprint}

\title{Calibrating Moral Intuitions}
\msid{MTE-08-b}
\msfamily{MTE-08}
\msdiscipline{Metaethics}
\mstype{Companion manuscript}
\mskeywords{moral intuition; reliability; calibration; framing effects; considered judgments}

\begin{document}
\maketitle

\begin{abstract}
MTE-08-a argues that moral intuition is a fallible but generally reliable faculty, comparable to perception under imperfect conditions, and that its reliability is a contingent matter, established by the kind of evidence that establishes the reliability of perception. We set out that evidence. An instrument is calibrated by three tests: agreement among instruments in good conditions, stability of readings as conditions improve, and an explanation of its errors by identifiable departures from good conditions. Moral intuition passes all three, and the experimental literature on framing effects, often cited against it, supplies evidence for the third.
\end{abstract}

\section{Three tests}\label{s:tests}

On the view of MTE-08-a, a moral intuition is a reaction to a particular case (Definition~2.1 of MTE-08-a), moral facts are facts about idealized reactions (Assumption~2.2), and reactions agree more often with idealized ones as the conditions of their formation approach the ideal (Assumption~2.4). Reliability is therefore a function of conditions (Lemma~3.2 of MTE-08-a), and conditions are measured by the subject's information, her vivid appreciation of what is at stake for those affected, and the equality of her attention to their interests (Definition~2.3 of MTE-08-a). The reliability of an instrument whose outputs cannot be checked against an independent reading is established by three tests, which are the ones used for perception before there were independent instruments.

\begin{definition}[Calibration evidence]\label{d:cal}
A faculty whose reliability varies with conditions is \emph{calibrated} by (i) \emph{agreement}: subjects whose conditions are good agree in their outputs; (ii) \emph{stability}: as a subject's conditions improve, her outputs change less and less; (iii) \emph{explanation of error}: where outputs disagree or change, the difference is explained by a departure from good conditions in at least one of them.
\end{definition}

\begin{proposition}\label{p:cal}
If moral intuition passes the three tests, then the hypothesis of MTE-08-a, that intuitions formed in good enough conditions agree with idealized reactions, is better supported than the hypothesis that intuitions are unreliable in good conditions.
\end{proposition}

\begin{proof}
On the first hypothesis, together with Assumption~2.4 of MTE-08-a, (i)--(iii) are predicted: idealized reactions converge, so reactions close to them agree; approaching them changes less as one gets closer; and disagreement arises only where some reaction is far from them. On the second hypothesis, agreement in good conditions would be coincidence, stability unexplained, and the correlation of error with departures from good conditions accidental.
\end{proof}

\section{The evidence}\label{s:ev}

\emph{Agreement.} \citet{Rawls1951} characterizes considered judgments as those formed by competent judges who know the relevant facts, are impartial, and are not affected by the outcome; these are the good conditions of MTE-08-a. In the core cases, the gratuitous infliction of suffering, the betrayal of trust for gain, the punishment of the known innocent, such judgments agree across traditions, and Ross's list of prima facie duties records a convergence that has not been seriously disturbed \citep{Ross1930}.

\emph{Stability.} Intuitions that change when subjects learn more or imagine more vividly change in a uniform direction. The expansion of concern to those previously excluded follows the acquisition of information about them and of imaginative contact with their situation, and it is not reversed by further information. This is the stability that (ii) requires.

\emph{Explanation of error.} \citet{SinnottArmstrong2008} argues from framing and order effects that moral intuitions are unreliable and need inferential confirmation. The effects are real, and they persist even among professional philosophers \citep{SchwitzgebelCushman2012}. But they appear where cases are presented in ways that distort the information or attention of the subject, and they are produced by just such departures from good conditions as Definition~2.3 of MTE-08-a counts: a description that renders one outcome more vivid than another is a departure in the vividness of appreciation; the order of cases shifts what is salient. That error correlates with these departures is evidence for test (iii), as the discovery of illusions under specific lighting confirms rather than undermines perception \citep[ch.~5]{Huemer2005}.

\begin{objection}
The tests are circular: good conditions are identified as those in which intuitions agree.
\end{objection}

\begin{reply}
They are not so identified. Information, vividness and impartiality are specified independently of agreement (Definition~2.3 of MTE-08-a), and each can be assessed without knowing what any subject concludes. The tests then check whether intuitions formed in conditions independently judged good agree, as perception is checked against conditions of lighting independently judged good.
\end{reply}

\section{Conclusion}

The reliability of moral intuition is not established by the concepts, but it is established. Agreement in good conditions, stability as conditions improve, and the explanation of errors by departures from good conditions are the marks of a calibrated faculty, and moral intuition has them.

\begin{thebibliography}{99}\small
\bibitem[Huemer(2005)]{Huemer2005} Huemer, M. (2005). \emph{Ethical Intuitionism}. Palgrave Macmillan.
\bibitem[Rawls(1951)]{Rawls1951} Rawls, J. (1951). Outline of a decision procedure for ethics. \emph{The Philosophical Review}, 60(2), 177--197.
\bibitem[Ross(1930)]{Ross1930} Ross, W.~D. (1930). \emph{The Right and the Good}. Oxford University Press.
\bibitem[Schwitzgebel and Cushman(2012)]{SchwitzgebelCushman2012} Schwitzgebel, E. and Cushman, F. (2012). Expertise in moral reasoning? Order effects on moral judgment in professional philosophers and non-philosophers. \emph{Mind \& Language}, 27(2), 135--153.
\bibitem[Sinnott-Armstrong(2008)]{SinnottArmstrong2008} Sinnott-Armstrong, W. (2008). Framing moral intuitions. In W. Sinnott-Armstrong (ed.), \emph{Moral Psychology, Volume 2: The Cognitive Science of Morality}. MIT Press.
\end{thebibliography}

\end{document}
